\documentclass[11pt]{article}
\usepackage[a4paper,margin=1in]{geometry}
\usepackage{amsmath,amssymb,amsthm,booktabs,enumitem,mathtools,microtype}
\usepackage[T1]{fontenc}
\usepackage{lmodern}
\usepackage[hidelinks]{hyperref}

\newtheorem{theorem}{Theorem}[section]
\newtheorem{proposition}[theorem]{Proposition}
\newtheorem{lemma}[theorem]{Lemma}
\newtheorem{corollary}[theorem]{Corollary}
\theoremstyle{definition}
\newtheorem{definition}[theorem]{Definition}
\theoremstyle{remark}
\newtheorem{remark}[theorem]{Remark}

\newcommand{\mob}{\mu}
\newcommand{\Ecal}{\mathcal E}
\newcommand{\Z}{\mathbb Z}
\newcommand{\kap}{\kappa_2}

\hypersetup{
  pdftitle={Prime-Square Tail Rate and a Quadratic Memory Remainder},
  pdfauthor={RH research program},
  pdfsubject={First-order square-clock gap rate in a fixed phasewise Mobius factor class},
  pdfkeywords={Mobius function, prime-square tail, Euler product, finite clocks, quantitative remainder}
}

\title{Prime-Square Tail Rate and a Quadratic Memory Remainder\\
for Phasewise Chowla-Free Memory}
\author{RH research program}
\date{August 7, 2026}

\begin{document}
\maketitle

\begin{abstract}
Within the fixed finite-clock, universally distance-two-safe phasewise
lag-two class whose interpolation coefficient $c_{11}(r)$ vanishes at every
phase, RH-379 identifies an all-clock supremum $B_\infty$ and RH-380 proves
that the square-clock values $G(q_y)$ increase strictly to it.  We determine
the first-order rate of that convergence.  Put
\[
 a_{j+1}=\frac1{p_{j+1}^2-1},\qquad
 T_y=\sum_{j\ge y}a_{j+1},
\]
where $3=p_1<p_2<\cdots$ are the odd primes, and let
\[
 e_m=\prod_{p\ \mathrm{odd}}\left(1-\frac m{p^2}\right),\qquad
 X_\infty=\frac{2e_4-4e_5+6e_6-8e_7+10e_8}{e_1}.
\]
The run formula gives $X_\infty\ge6e_8/e_1>0$.  A factorwise Euler-tail
comparison proves $|X_j-X_\infty|\le170T_j$.  The exact RH-379 product for
$H_{j+1}$ gives
\[
 0\le\frac4{\pi^2}-H_{j+1}
 \le\frac4{\pi^2}T_{j+1},
\]
and the normalized memory statistic satisfies $0\le M_j/A_j\le1$.
Telescoping the exact RH-380 increment and using two exact tail-sum
identities yields, for every $y\ge1$,
\[
 \left|B_\infty-G(q_y)-\frac{2X_\infty}{\pi^2}T_y\right|
 \le \frac{342}{\pi^2}T_y^2.
\]
Consequently
\[
 \frac{B_\infty-G(q_y)}{T_y}\longrightarrow
 \frac{2X_\infty}{\pi^2}>0.
\]
No prime number theorem is used.  The result supplies neither a second-order
coefficient nor a $p_y$-scale asymptotic, growing-clock theorem, adaptive
capacity limit, operator, trace formula, zero identification, or implication
for the Riemann hypothesis.
\end{abstract}

\noindent\textbf{Keywords:} M\"obius function; prime-square Euler tail;
square clocks; quantitative convergence; lag-two memory.

\section{Frozen class and exact predecessor inputs}

The theorem concerns exactly the class isolated in RH-379 and retained in
RH-380 \cite{RH379,RH380}.  We begin by fixing its quantifiers.

\begin{definition}[Fixed-clock phasewise class]
Fix an integer $q\ge1$ before taking $N\to\infty$.  For every
$r\in\Z/q\Z$, choose a table
\[
 f_r:\{-1,0,1\}^2\longrightarrow\{-1,+1\},\qquad
 \epsilon_n=f_{n\bmod q}\bigl(\mob(n-2),\mob(n)\bigr).
\]
The family is universally distance-two-safe if no ternary input can give
$\epsilon_n=\epsilon_{n+2}=+1$.  It is in the present class only when the
coefficient $c_{11}(r)$ of $\mob(n-2)\mob(n)$ in the RH-379 interpolation is
zero for every phase.  The exact fixed-clock optimum is denoted by $G(q)$.
\end{definition}

The phasewise condition is not cosmetic.  If $c_{11}(r)\ne0$, the limiting
expansion contains a phase-weighted shift-two M\"obius correlation that is
not controlled by the frozen sources.  Nothing below averages or cancels
those missing terms.

Let $3=p_1<p_2<\cdots$ be the odd primes and define the square-clock data
\begin{equation}
 P_y=\prod_{i\le y}p_i^2,\qquad q_y=4P_y,\qquad
 A_y=\prod_{i\le y}(p_i^2-1),\qquad
 D_y=\prod_{i\le y}(p_i^2-2).
 \label{eq:square-data}
\end{equation}
RH-374 associates a cyclic odd squarefree-support word to $P_y$.  Its
positive runs have lengths at most eight; write $R_\ell^{(y)}$ for the
number of runs of length $\ell$ \cite{RH374}.  With
\begin{equation}
 E_m^{(y)}=\prod_{i\le y}\left(1-\frac m{p_i^2}\right),
 \label{eq:finite-euler}
\end{equation}
the exact run formula is
\begin{equation}
 R_\ell^{(y)}=
 \begin{cases}
 P_y\bigl(E_\ell^{(y)}-2E_{\ell+1}^{(y)}+E_{\ell+2}^{(y)}\bigr),
      &1\le\ell\le7,\\
 P_yE_8^{(y)},&\ell=8.
 \end{cases}
 \label{eq:run-formula}
\end{equation}
Set
\begin{align}
 \Ecal_y&=R_2^{(y)}+R_4^{(y)}+R_6^{(y)}+R_8^{(y)},
 &L_y&=2R_2^{(y)}+4R_4^{(y)}+6R_6^{(y)}+8R_8^{(y)},
 \label{eq:EL}\\
 M_y&=2R_3^{(y)}+4R_5^{(y)}+6R_7^{(y)}.
 \label{eq:M}
\end{align}
RH-380 proves the exact increment
\begin{equation}
\begin{aligned}
G(q_{j+1})-G(q_j)
={}&\frac{2(L_j-2\Ecal_j)}{\pi^2A_j(p_{j+1}^2-1)}\\
&+\frac{M_j}{A_j(p_{j+1}^2-1)}
\left(\frac4{\pi^2}-H_{j+1}\right),
\end{aligned}
\label{eq:rh380-increment}
\end{equation}
where $H_j=\kap A_j/D_j$, and it records
\begin{equation}
 G(q_j)\longrightarrow B_\infty.
 \label{eq:G-limit}
\end{equation}
The limit in \eqref{eq:G-limit} is the RH-379 all-clock supremum.  It is a
cofinal limit after the fixed-clock $N$-limits, not a growing choice
$q=q(N)$.

\section{Prime-square tails and the normalized run statistic}

For $j\ge1$, put
\begin{equation}
 a_{j+1}=\frac1{p_{j+1}^2-1},\qquad
 T_j=\sum_{k\ge j}a_{k+1}.
 \label{eq:T}
\end{equation}
The notation deliberately keeps the successor index visible: the increment
from $q_j$ to $q_{j+1}$ carries $a_{j+1}$.

For $1\le m\le8$, define
\begin{equation}
 e_m=\prod_{p\ \mathrm{odd}}\left(1-\frac m{p^2}\right),\qquad
 U_m^{(j)}=\frac{E_m^{(j)}}{E_1^{(j)}},\qquad
 u_m=\frac{e_m}{e_1}.
 \label{eq:ratios}
\end{equation}
All factors in these products are positive.  Their convergence and
positivity follow elementarily from convergence of $\sum_{n\ge3}n^{-2}$.
The RH-374 normalization is $e_1=8/\pi^2$ \cite{RH374}.

\begin{definition}[Normalized deterministic numerator]
Define
\begin{equation}
 X_j=\frac{L_j-2\Ecal_j}{A_j},
 \qquad
 X_\infty=2u_4-4u_5+6u_6-8u_7+10u_8.
 \label{eq:X-def}
\end{equation}
Equivalently,
\[
 X_\infty=
 \frac{2e_4-4e_5+6e_6-8e_7+10e_8}{e_1}.
\]
\end{definition}

\begin{lemma}[Exact finite Euler form]
\label{lem:X-euler}
For every $j\ge1$,
\begin{equation}
 X_j=2U_4^{(j)}-4U_5^{(j)}+6U_6^{(j)}
      -8U_7^{(j)}+10U_8^{(j)}.
 \label{eq:X-euler}
\end{equation}
Moreover,
\begin{equation}
 X_\infty\ge\frac{6e_8}{e_1}>0.
 \label{eq:X-positive}
\end{equation}
\end{lemma}

\begin{proof}
The run identity in RH-380 is
\[
 L_j-2\Ecal_j=2R_4^{(j)}+4R_6^{(j)}+6R_8^{(j)}.
\]
Since $A_j=P_jE_1^{(j)}$, substituting \eqref{eq:run-formula} gives
\begin{align*}
X_j
&=2\frac{E_4^{(j)}-2E_5^{(j)}+E_6^{(j)}}{E_1^{(j)}}
 +4\frac{E_6^{(j)}-2E_7^{(j)}+E_8^{(j)}}{E_1^{(j)}}
 +6\frac{E_8^{(j)}}{E_1^{(j)}}\\
&=2U_4^{(j)}-4U_5^{(j)}+6U_6^{(j)}
 -8U_7^{(j)}+10U_8^{(j)}.
\end{align*}
Passing to convergent Euler products proves the formula for $X_\infty$.
The two finite second differences in the first line are normalized run
counts and hence nonnegative.  Their limits remain nonnegative.  Dropping
them leaves $6e_8/e_1$, which is strictly positive.
\end{proof}

\begin{lemma}[Factorwise tail bound]
\label{lem:X-bound}
For every $j\ge1$,
\begin{equation}
                 \boxed{|X_j-X_\infty|\le170T_j.}
 \label{eq:X-bound}
\end{equation}
\end{lemma}

\begin{proof}
For $m\in\{4,5,6,7,8\}$, division of the Euler factors gives the exact
tail identity
\begin{equation}
 \frac{u_m}{U_m^{(j)}}
 =\prod_{p>p_j}\frac{1-m/p^2}{1-1/p^2}
 =\prod_{p>p_j}\left(1-\frac{m-1}{p^2-1}\right).
 \label{eq:ratio-tail}
\end{equation}
For numbers $0\le v_i\le1$, the elementary product union bound
$0\le1-\prod_i(1-v_i)\le\sum_i v_i$ follows first for finite products and
then by monotone passage to the limit.  Also $0<U_m^{(j)}\le1$.  Therefore
\begin{equation}
 0\le U_m^{(j)}-u_m
 \le (m-1)T_j.
 \label{eq:one-ratio-bound}
\end{equation}
Apply the triangle inequality to \eqref{eq:X-euler}.  The coefficient
ledger is
\[
 2\cdot3+4\cdot4+6\cdot5+8\cdot6+10\cdot7
 =6+16+30+48+70=170,
\]
which proves \eqref{eq:X-bound}.
\end{proof}

\section{The \texorpdfstring{$H$}{H} tail and the memory statistic}

\begin{lemma}[Exact $H$ product and its tail loss]
\label{lem:H-tail}
For every $j\ge1$,
\begin{equation}
 \frac{H_{j+1}}{4/\pi^2}
 =\prod_{p>p_{j+1}}\left(1-\frac1{p^2-1}\right),
 \label{eq:H-product}
\end{equation}
and hence
\begin{equation}
 0\le\frac4{\pi^2}-H_{j+1}
 \le\frac4{\pi^2}T_{j+1}.
 \label{eq:H-bound}
\end{equation}
\end{lemma}

\begin{proof}
The exact RH-379 product is $H_j=\kap A_j/D_j$, where
$\kap=\prod_p(1-2/p^2)$.  Separating the prime $2$ and cancelling the
finite odd-prime factors gives
\[
 H_{j+1}=\frac12 E_1^{(j+1)}
 \prod_{p>p_{j+1}}\left(1-\frac2{p^2}\right).
\]
Because $4/\pi^2=e_1/2$, division yields
\[
 \frac{H_{j+1}}{4/\pi^2}
 =\prod_{p>p_{j+1}}
 \frac{1-2/p^2}{1-1/p^2}
 =\prod_{p>p_{j+1}}\left(1-\frac1{p^2-1}\right).
\]
The same product union bound used in Lemma~\ref{lem:X-bound} now gives
\eqref{eq:H-bound}.
\end{proof}

\begin{lemma}[Quadratic-memory normalization]
\label{lem:M-bound}
For every $j\ge1$,
\begin{equation}
                         0\le\frac{M_j}{A_j}\le1.
 \label{eq:M-bound}
\end{equation}
\end{lemma}

\begin{proof}
Nonnegativity is immediate from \eqref{eq:M}.  An odd run of length
$\ell$ contributes $\ell-1\le\ell$ to $M_j$.  Thus $M_j$ is at most the
number of positive sites lying in odd runs, which is at most the total
number of positive sites.  The latter is exactly
$P_jE_1^{(j)}=A_j$.  This proves \eqref{eq:M-bound}.
\end{proof}

The word ``quadratic'' in the title refers to the contribution of this
memory statistic to the final $T_y^2$ remainder.  It does not announce a
quadratic coefficient expansion.

\section{Two exact tail-sum identities}

\begin{lemma}[Tail algebra]
\label{lem:tail-algebra}
For every $y\ge1$,
\begin{align}
 \sum_{j\ge y}a_{j+1}T_j
 &=\frac12\left(T_y^2+\sum_{j\ge y}a_{j+1}^2\right)
 \le T_y^2,
 \label{eq:tail-current}\\
 \sum_{j\ge y}a_{j+1}T_{j+1}
 &=\frac12\left(T_y^2-\sum_{j\ge y}a_{j+1}^2\right)
 \le\frac12T_y^2.
 \label{eq:tail-next}
\end{align}
\end{lemma}

\begin{proof}
Since $T_{j+1}=T_j-a_{j+1}$,
\begin{align*}
 a_{j+1}T_j
 &=\frac12\bigl(T_j^2-T_{j+1}^2+a_{j+1}^2\bigr),\\
 a_{j+1}T_{j+1}
 &=\frac12\bigl(T_j^2-T_{j+1}^2-a_{j+1}^2\bigr).
\end{align*}
Sum these identities from $j=y$ to a finite endpoint and then let that
endpoint tend to infinity.  The tails tend to zero because
\begin{equation}
 0<T_y\le\sum_{n>p_y}\frac1{n^2-1}
 =\frac12\left(\frac1{p_y}+\frac1{p_y+1}\right)\longrightarrow0.
 \label{eq:T-zero}
\end{equation}
Here positivity holds for every finite $y$ because there are infinitely
many odd primes.  Finally,
$\sum a_{j+1}^2\le(\sum a_{j+1})^2=T_y^2$, which gives both displayed
inequalities.  No estimate for the distribution of primes is involved.
\end{proof}

\section{First-order gap rate}

\begin{proposition}[Exact infinite increment sum]
\label{prop:infinite-sum}
For every $y\ge1$,
\begin{equation}
\begin{aligned}
B_\infty-G(q_y)
=\sum_{j\ge y}\biggl[
 &\frac{2X_j}{\pi^2}a_{j+1}\\
 &+\frac{M_j}{A_j}a_{j+1}
 \left(\frac4{\pi^2}-H_{j+1}\right)
\biggr].
\end{aligned}
\label{eq:infinite-sum}
\end{equation}
\end{proposition}

\begin{proof}
For a finite $K>y$, telescope the exact RH-380 increment
\eqref{eq:rh380-increment} from $j=y$ through $K-1$ and use
$X_j=(L_j-2\Ecal_j)/A_j$ and
$a_{j+1}=(p_{j+1}^2-1)^{-1}$.  Then let $K\to\infty$ and invoke the
already proved cofinal limit \eqref{eq:G-limit}.  This is a limit of finite
telescoping identities after all fixed-clock $N$-limits have been taken;
there is no exchange of an $N$-limit with a $j$-limit.
\end{proof}

\begin{theorem}[Prime-square tail rate]
\label{thm:rate}
For every $y\ge1$,
\begin{equation}
 \boxed{
 \left|B_\infty-G(q_y)-\frac{2X_\infty}{\pi^2}T_y\right|
 \le\frac{342}{\pi^2}T_y^2.}
 \label{eq:rate-bound}
\end{equation}
Consequently,
\begin{equation}
 \boxed{
 \frac{B_\infty-G(q_y)}{T_y}
 \longrightarrow\frac{2X_\infty}{\pi^2}>0.}
 \label{eq:ratio-limit}
\end{equation}
\end{theorem}

\begin{proof}
Subtract $(2X_\infty/\pi^2)T_y$ from
\eqref{eq:infinite-sum}.  Since $T_y=\sum_{j\ge y}a_{j+1}$, the first
remainder is bounded using Lemma~\ref{lem:X-bound} and
\eqref{eq:tail-current}:
\begin{align}
\left|\sum_{j\ge y}\frac{2(X_j-X_\infty)}{\pi^2}a_{j+1}\right|
&\le\frac{340}{\pi^2}\sum_{j\ge y}a_{j+1}T_j
\le\frac{340}{\pi^2}T_y^2.
\label{eq:X-remainder}
\end{align}
The memory summand is nonnegative.  Lemmas~\ref{lem:H-tail} and
\ref{lem:M-bound}, followed by \eqref{eq:tail-next}, give
\begin{align}
0&\le\sum_{j\ge y}\frac{M_j}{A_j}a_{j+1}
 \left(\frac4{\pi^2}-H_{j+1}\right)\\
&\le\frac4{\pi^2}\sum_{j\ge y}a_{j+1}T_{j+1}
\le\frac2{\pi^2}T_y^2.
\label{eq:M-remainder}
\end{align}
Adding \eqref{eq:X-remainder} and \eqref{eq:M-remainder} proves
\eqref{eq:rate-bound}; the ledger is $340+2=342$.

For each finite $y$, $T_y>0$ by \eqref{eq:T-zero}.  Divide
\eqref{eq:rate-bound} by $T_y$, then use $T_y\to0$ from the same elementary
integer-square comparison.  Equation \eqref{eq:ratio-limit} follows, and
its limit is positive by \eqref{eq:X-positive}.
\end{proof}

\begin{remark}[No prime number theorem]
The scale in Theorem~\ref{thm:rate} is the intrinsic prime-square tail
$T_y$.  We never replace it by a function of $p_y$.  The only tail
comparison is \eqref{eq:T-zero}, obtained by telescoping the integer series
$(n^2-1)^{-1}$.  Thus the proof uses neither the prime number theorem nor a
weaker prime-counting asymptotic.
\end{remark}

\section{Exact artifact and finite diagnostics}

The standard-library artifact has two distinct layers.  First, exact
\texttt{Fraction} arithmetic regenerates the run statistics, normalized
Euler ratios, RH-380 increments, and finite versions of both identities in
Lemma~\ref{lem:tail-algebra}.  The independently frozen six-row fixture has
canonical SHA-256
\begin{center}
\texttt{d55fd48071eb5b88c054f3d34329f274f792f2bbd859b4ab98e31b5b7020beb8}.
\end{center}
Second, directed decimal arithmetic at precision $60$ is converted to
exact rational endpoints.  It enumerates the $9592$ primes through
$100000$ and bounds every omitted prime by the exact integer tail
\begin{equation}
 \sum_{n>100000}\frac1{n^2-1}
 =\frac{200001}{20000200000}.
 \label{eq:numeric-tail}
\end{equation}
All comparisons fail closed if outward upper and lower endpoints overlap in
the wrong order.  The rows $y=1,2,3,5,10,25$ reproduce the $170$ and $342$
bounds; they do not prove the all-$y$ theorem by fitting.

For orientation only, the directed product layer encloses
\[
 0.46967971561917<X_\infty<0.46993774873491.
\]
These decimals are a finite diagnostic.  Positivity in the theorem is the
exact run argument \eqref{eq:X-positive}, not this numerical enclosure.

The release builder locks exactly $25$ immutable inputs: seven RH-374
files, eight RH-379 files, eight RH-380 files, and two RH-MVP2 archive
summaries.  It checks both the live SHA-256 and byte identity with each
declared release commit.  Mutable project instructions and handoff state
are deliberately excluded.  The result ledger has a recursively closed
Draft 2020-12 schema, and archive replay rejects missing members, unsafe
paths, duplicate JSON keys, hash drift, release-commit rebinding, or a
semantic PDF that differs from \texttt{main.pdf}.

Reproduction from this directory is
\begin{verbatim}
make result
make test
make pdf
make archive
\end{verbatim}

\section{Scope, limitations, and next edge}

The theorem closes one precise question: it identifies the first-order
prime-square-tail scale of the square-clock nonattainment gap and proves an
explicit quadratic remainder inside the frozen RH-379 class.  Its limits
are equally precise.

\begin{itemize}[leftmargin=2em]
\item The theorem is phasewise $c_{11}=0$.  It does not cover unrestricted
      lag-two tables or assume the missing phase-weighted shift-two
      correlation cancellation.
\item The constant $342$ is explicit and uniform for $y\ge1$, but no
      optimality is claimed for either $170$ or $342$.
\item No exact second-order coefficient is identified.  In particular,
      the theorem does not collapse all quadratic information to one
      multiple of $T_y^2$.
\item No asymptotic for $T_y$ in terms of $p_y$, no growing clock $q(N)$,
      and no adaptive-capacity convergence is proved.
\item Finite product rows and decimal enclosures are reproduction only;
      they are not finite-order regression promoted to an all-order claim.
\item No canonical intrinsic dynamical spectral determinant,
      time-oriented scattering completion, self-adjoint generator,
      von Mangoldt prime-power trace identity, or completed-zeta divisor
      equality is constructed.  Gates A--E remain false/open
      \cite{RHMVP2}.
\item There is no Hilbert--P\'olya operator, no identification of Riemann
      zeros, and no implication for the Riemann hypothesis.
\end{itemize}

The immediate within-class next edge is a genuinely second-order analysis
of the Euler tails and of $M_j/A_j$.  Such a successor must retain every
independent quadratic tail scale that survives the exact algebra and prove
an all-order remainder; this paper does not state its coefficient.  The
first class-enlargement trigger remains a theorem for the required
phase-weighted shift-two M\"obius correlations.

\section*{Declarations}

\paragraph{Data and code availability.}
All theorem inputs, exact source code, tests, closed result schema, source
locks, and archive manifests are contained in the repository paper
directory.  No private dataset is used.

\paragraph{Author contributions.}
The RH research program performed conceptualization, formal analysis,
software construction, verification, artifact curation, adversarial review,
and manuscript preparation.

\paragraph{Funding.}
No external funding is declared.

\paragraph{Competing interests.}
No competing interests are declared.

\paragraph{Ethics and human participants.}
This mathematical and computational study uses no human participants,
animals, or personal data; ethics approval and consent are not applicable.

\paragraph{AI assistance disclosure.}
AI-assisted tools supported source triage, symbolic checking, deterministic
artifact generation, drafting, and adversarial review.  Every released
claim is constrained by immutable repository sources and reproducible
verification artifacts; responsibility for the scoped mathematical
statements remains with the research program.

\bibliographystyle{plain}
\bibliography{references}

\end{document}
